% 第 41 章　弯曲时空的邀请函（选读）
\chapter{弯曲时空的邀请函（选读）}

第 38 章到第 40 章里，相对论的舞台是“平直”的：没有引力，惯性系铺满整个时空，钟和尺的规矩由光速不变一条原理说了算。可这套理论里有一个巨大的空缺——引力。牛顿的万有引力定律在低速世界精确得近乎完美，却始终没能和光速上限握手言和。本章是爱因斯坦给引力写的一封“邀请函”：他请引力放弃“力”的身份，改头换面成为时空本身的几何。这一章只画概念地图，不推导广义相对论的场方程；我们会说清楚哪些已经是千锤百炼的实验事实，哪些是目前最好的数学模型，哪些解释仍悬而未决。

\step{反常识开场}

\begin{mainline}
掏出手机，打开地图，看着那个代表你的蓝点稳稳地落在马路上。这个再普通不过的动作，背后藏着一笔每天 38 微秒的账。全球定位系统的卫星在两万千米高空绕地球飞行，每颗卫星都带着一台原子钟。地面控制中心的工程师发现：如果不给这些钟做任何修正，卫星钟每天都会比地面钟快约 38 微秒。38 微秒听起来微不足道，但定位靠的正是“无线电信号飞了多久”，而无线电以光速传播，38 微秒乘上光速就是约 11 千米——也就是说，不修正这笔账，你的导航每天会多漂出约 10 千米，不出一周，蓝点就会把你标进隔壁城市。工程师的解决办法很朴素：卫星上天之前，预先把原子钟的节拍调慢一点点，让它在轨道上刚好和地面同步。这不是质量瑕疵，而是按计划补偿一个真实存在的物理效应：高处的钟，天生比低处的钟走得快。

再看一个画面。空间站里的宇航员松开一支笔，笔就静静地飘在半空。很多人会解释：“太空里没有重力嘛。”可是算一下就露馅了：空间站离地只有四百千米上下，那里的地球引力约是地面的九成——如果宇航员站在一杆（想象中的）足够高的秤上，秤会忠实地显示出接近地面体重的读数。引力明明还在，人却完全感觉不到它。引力这种“可以凭空消失”的脾气，是电磁力、摩擦力从来没有的。把你关进一间没有窗的电梯，切断钢缆，你会和电梯一起自由下落，舱内所有物体一起飘浮——在舱内做任何实验，你都查不出引力的存在。

一台必须靠“调慢时钟”才能工作的导航系统，一种“一松手就消失”的力。这两条线索指向同一个惊人的结论：引力也许根本就不是一种力。
\end{mainline}

\step{先猜一猜}

\begin{mainline}
先把你此刻的直觉摆到桌面上，凭第一感觉判断下面四个说法：

第一，光没有质量，所以引力拉不动光，光永远走直线。

第二，只要不被任何东西支撑，在哪里做实验都一样：山顶的钟和海平面的钟，走得一样快。

第三，把你关在一间没有窗户、隔音良好的电梯里，只要电梯足够平稳，你总能做一个实验，分辨出自己是“停在地面上的电梯里”，还是“在没有引力的深空里被火箭以 \(9.8\ \mathrm{m/s^{2}}\) 推着加速”。

第四，引力是一种真实的力，就像电磁力一样，从太阳出发，“作用”在地球上，拉着地球绕圈。

请记下你的答案。读完本章再回来对照：第一条和第二条会被实验直接否决；第三条恰恰是爱因斯坦最关键的思想实验，答案出乎所有人意料；第四条则是本章要拆掉的旧框架——拆掉之后，地球绕太阳这件事会有一个全新的、更美的说法。
\end{mainline}

\step{动手实验}

\begin{experiment}[用手机加速度计“抓住”引力]
你的手机里有一枚微型加速度计，平时负责横竖屏切换和计步。下载任意一款能显示传感器读数的应用（搜索“加速度计”即可找到一大批免费的），然后做三个动作。

第一步：把手机平放在桌面上。你会看到读数稳定在约 \(9.8\ \mathrm{m/s^{2}}\)，方向向上。手机明明静止不动，加速度计却报告一个向上的加速度——这正是本章的题眼：加速度计测的不是“动不动”，而是“有没有被支撑”。桌面一直在向上推着手机，传感器感受到的是这份支撑。

第二步：在床上垫好枕头，把手机轻轻抛起再接住（务必低空、务必接住）。在离手的那一瞬间，读数会跳到接近零。手机在真正“加速下落”，传感器反而报告“没有加速度”。引力的效应从读数上消失了——和空间站里的宇航员一模一样。

第三步：拿着手机走进电梯。启动上升的瞬间，读数短暂变大；减速停靠时，读数短暂变小。如果此时你脚下踩着一台体重秤，秤的读数会同步地变重、变轻。电梯用加速“伪造”了引力的大小变化。

这三个动作合在一起，就是爱因斯坦思想实验的掌上版：支撑能被感觉到，引力本身感觉不到；加速和引力在你身体的感觉里可以互相冒充。下次坐过山车或跳楼机，如果在安全范围内让手机固定在身前，你会在俯冲的那一秒看到读数再次跌向零——那一秒里，你和舱内飘浮的宇航员遵循的是同一条物理。
\end{experiment}

\step{旧模型遇到麻烦}

\begin{mainline}
牛顿的万有引力定律用了两百多年，精确到能预报海王星的存在。但它身上有四道裂缝，每一道都补不上。

第一道裂缝：引力是瞬时的。按牛顿定律，太阳此刻消失，地球此刻就飞出轨道——引力的消息传得无限快。可第 38 章已经立下铁律：任何携带能量和信息的影响都不能超过光速。如果引力消息瞬时到达，就可以用它造出超光速电报，因果律当场崩溃。牛顿本人对此也不安，他称超距作用“荒谬至极”，只是他找不到替代方案。

第二道裂缝：一个精确得诡异的巧合。牛顿力学里有两种质量：惯性质量决定“推起来有多费劲”，引力质量决定“被引力拉得有多狠”。它们概念上毫无关系，却严格相等——于是所有物体自由下落的加速度完全相同。这不是理论预言，是天上掉下来的巧合。从伽利略的斜面到 19 世纪厄缶的扭秤，再到 2022 年法国“显微镜号”卫星把两种材料的测试块放进轨道比较下落，实验把这个巧合压到了 \(10^{-15}\) 的精度：两个物体下落加速度的相对差别小于千万亿分之一。一个理论越精确地依赖巧合，越说明巧合背后藏着没被发现的原因。

第三道裂缝：水星的轨道差了一点点。牛顿力学能解释水星近日点进动的绝大部分——每世纪约 5600 角秒里，其他行星的拉扯、岁差等因素解释了 5557 角秒，剩下每世纪 43 角秒，无论如何对不上账。43 角秒大约是三千米外一根头发丝的宽度，19 世纪的天文学家却宁可怀疑存在一颗看不见的“火神星”，也不肯怀疑牛顿定律。火神星始终没有找到。

第四道裂缝：光。按牛顿的思路，引力按质量拉东西，光子没有静质量，光理应笔直穿过引力场。可是你在动手实验里已经看到了：加速和引力分不清。那么在加速的电梯里，光怎么走？这个问题牛顿框架根本回答不了——它的词典里没有“引力影响时空本身”这个条目。

四道裂缝指向同一个方向：引力不像一种普通的力，倒像是舞台本身的性质。
\end{mainline}

\step{发明新概念}

\begin{mainline}
1907 年，爱因斯坦在专利局里想到了他自称“一生中最快乐的思想”。设想你待在一间没有窗户的电梯里，脚下踩着体重秤。情形一：电梯停在地面，秤显示你的正常体重。情形二：电梯在没有任何引力的深空，被火箭以 \(9.8\ \mathrm{m/s^{2}}\) 的加速度向上推——地板加速撞向你的脚，秤同样显示正常体重。你在舱内做力学实验：松手的球以 \(9.8\ \mathrm{m/s^{2}}\)“落”向地板，单摆照常摆动。情形一和情形二，舱内任何实验都无法区分。反过来：情形三，电梯的钢缆断了，舱和人一起自由下落，舱内一切飘浮；情形四，电梯在深空惯性漂浮，远离一切星体，舱内一切同样飘浮。情形三和情形四也无法区分。

这就是等效原理：在一个足够小的局部区域里，均匀引力场的效果和一个加速参考系的效果不可区分。注意“局部”两个字——它后面会成为重要的边界。

先别急着往下走，回头看一眼：旧模型里那道最诡异的裂缝，已经被这个原理悄悄焊上了。为什么惯性质量和引力质量严格相等、为什么羽毛和铁球在真空里落得一样快？在“引力是一种力”的框架里，这是 \(10^{-15}\) 精度的无巧不成书；在等效原理的框架里，这根本不是一个需要解释的性质——自由下落的物体只是沿着时空里自己的路走，路径由时空的几何决定，几何才不管你是铁还是羽毛。伽利略在斜面上发现的那个“巧合”，在这里升格为理论的地基。

等效原理立刻给出两个可检验的预言。第一个是光会弯。设想电梯在深空向上加速，从舱壁水平射入一束激光。光飞渡船舱的短时间里，舱已经向上挪了一点点，于是光打在对面舱壁上的位置，比入射点低了一截——在舱内的人看来，光线向下弯成一条抛物线。既然加速舱和引力场不可区分，那么在地面上水平射出的光，也该同样向下弯。引力不是“拉住”了有质量的光子，而是连光走的“直线路径”本身都变了样。

第二个预言更安静，却更深刻：引力会改变钟的快慢。还是那部加速的电梯，光从地板射向天花板。光飞行的一小段时间里，天花板在加速远离，收到光时它的速度已经比发光时快了一截——用你熟悉的救护车多普勒效应类推，天花板收到的光频率会变低。既然加速舱等效于引力场，那么在引力场中，光从低处爬到高处，频率也会降低，波长变长——这叫引力红移。现在把逻辑再推一步：光的每一次振动就是一次“滴答”，高处收到的滴答次数少，唯一的自洽解释是低处的钟本身走得慢。不是钟坏了，是时间本身在不同高度流速不同。1960 年，庞德和雷布卡在哈佛大学一座 22.5 米高的塔里，用极其敏锐的核共振技术测到了塔顶与塔底光子频率的相对差别：约 \(2.5\times10^{-15}\)，与预言相符。1971 年，哈菲勒和基廷把原子钟带上民航客机环球飞行，回来后和地面的钟对表，差别与理论相符，量级是几百纳秒。时间随高度变化，早已不是思想实验，是测过一遍又一遍的事实。

到这里，爱因斯坦完成了一次惊险的跳跃。既然引力在局部可以被加速“消掉”，既然连钟的快慢和光的路径都随引力改变，那么引力也许根本不是物体之间传递的一种力，而是时空这张“地图”本身的形状。质量压弯了时空，物体只是在弯曲的时空里走它所能走的最直的路——就像飞机沿地球大圆航线飞行，飞行员从未转弯，航迹在平面地图上却弯成弧线。物理学家惠勒把整套理论浓缩成一句顺口溜：物质告诉时空如何弯曲，时空告诉物质如何运动。1915 年，爱因斯坦把这句顺口溜写成了精确的方程，广义相对论诞生。它顺手解释了水星那 43 角秒——不多不少，正好是弯曲时空的修正。
\end{mainline}

\step{一句话模型}

\begin{oneline}
引力不是一种从远处拉你的力，而是时空本身被物质压弯后的形状：自由下落的物体走的是弯曲时空里“最直的路”，而时间的快慢也是这形状的一部分——离大质量物体越近，钟走得越慢。你在地面感到的“体重”，其实是大地顶着你不让你沿最直的路走下去。
\end{oneline}

\step{数学翻译器}

\begin{mainline}
本章遵守约定，不碰场方程的推导，但有三条“可计算的直觉”值得翻译出来，它们全都可以用初中代数检查。

第一条：弱引力场里的钟差公式。两只钟的引力势相差 \(\Delta\Phi\)（在地球表面附近，高度差 \(h\) 对应 \(\Delta\Phi=gh\)），它们快慢的相对差别是
\[
\frac{\Delta f}{f}=\frac{\Delta\Phi}{c^{2}}
\]
回答四个问题。描述什么：两只钟计时快慢的相对差别。为什么这样写：分母必须是能量量纲，自然界能用的标尺只有 \(c^{2}\)，所以差别必然被压得极小——这正是“时间弯曲”日常完全感觉不到的定量原因。何时成立：弱引力场（地球、太阳附近都没问题），钟近似静止或速度修正另行计算。何时失效：黑洞视界附近、中子星内部这类强场，要用完整的广义相对论重新算。

立刻做特殊检查。\(h=0\)：钟差为零，同一高度的钟同步，合理。\(c\) 换成无限大：差别消失——牛顿世界里时间是绝对的，公式自动退回旧理论，合理。方向检查：把两只钟上下对调，\(\Delta\Phi\) 变号，快慢关系反转，合理。数量级检查：珠穆朗玛峰顶比海平面高约 9 千米，代入得 \(\Delta f/f\approx10^{-12}\)，一年快约 30 微秒——真实存在，小得只有原子钟在乎。
\end{mainline}

\begin{mainline}
第二条：GPS 的账，用真实数字算一遍。地球的质量和引力常数的乘积 \(GM\approx4\times10^{14}\ \mathrm{m^{3}/s^{2}}\)，引力势写作 \(\Phi=-GM/r\)。地面半径约 \(6.4\times10^{6}\) 米，GPS 卫星轨道半径约 \(2.66\times10^{7}\) 米。钟差的引力部分：
\[
\frac{\Delta\Phi}{c^{2}}=\frac{GM}{c^{2}}\left(\frac{1}{R_{\text{地}}}-\frac{1}{r_{\text{卫}}}\right)\approx5.3\times10^{-10}
\]
乘上一天的 86400 秒，卫星钟每天因高度快约 46 微秒。但卫星还在以约每秒 3.9 千米的速度飞奔，第 39 章的时间膨胀要它慢：\(v^{2}/(2c^{2})\approx8\times10^{-11}\)，每天慢约 7 微秒。两笔账一合并：每天净快约 38 微秒。这 38 微秒乘上光速，就是开场里那 11 千米的漂移。注意一个关键事实：广义相对论那 46 微秒不但不小，反而是狭义相对论修正的六倍多——如果只会狭义相对论，GPS 照样修不对。你的手机每天免费使用着 1915 年的理论。
\end{mainline}

\begin{mainline}
第三条：光掠过太阳时弯多少。用等效原理加量纲估算，可以得到一个角秒量级的偏折；1915 年完整理论的精确预言是：贴着太阳表面掠过的星光偏折 \(1.75\) 角秒——约等于五千米外一枚硬币的张角。有趣的是，把光子当成被太阳吸引的“小颗粒”做牛顿式估算，得到的恰好是 \(0.87\) 角秒，只有真值的一半。差的那一半去哪了？弯曲的不只是时间，还有空间——等效原理的电梯只抓住了“时间弯曲”，完整理论里空间的几何也弯了，两半相加才对。这个“一半”是极好的警示：直觉能把你送到正确的大门，却未必送得进门。1919 年，爱丁顿带队趁日全食拍摄太阳附近的恒星位置，测到的位移与 \(1.75\) 角秒相符，爱因斯坦一夜之间举世闻名。现代用射电干涉仪重复这类测量，精度已达 \(10^{-4}\) 量级，理论与观测始终一致。
\end{mainline}

\begin{center}
\begin{tikzpicture}[scale=1.0,>=Stealth]
  % 加速电梯中的光束
  \draw[thick] (0,0) rectangle (3.2,4.6);
  \node at (1.6,-0.35) {加速上升的电梯舱};
  \draw[->,thick] (3.2,2.3) -- (4.2,2.3) node[right] {加速度 \(g\)};
  % 光束入射与弯曲路径
  \fill (0,3.6) circle (2pt);
  \draw[->,thick] (0,3.6) .. controls (1.6,3.35) and (2.6,2.9) .. (3.2,2.3);
  \fill (3.2,2.3) circle (2pt);
  % 水平参考线
  \draw[dashed] (0,3.6) -- (3.2,3.6);
  \draw[<->] (2.9,3.6) -- (2.9,2.3) node[midway,right] {偏低};
  \node at (1.2,4.1) {入射光束};
\end{tikzpicture}
\end{center}

\begin{mainline}
这幅图就是等效原理的全部工作方式：舱向上加速，光在舱内画出一道下弯的弧线；把“加速舱”换成“静止在引力场中的舱”，弧线原样保留。下一幅图把同一逻辑搬到太阳旁边：远处的恒星发出的光掠过太阳边缘，路径偏折一个微小角度，地球上的我们顺着光的来路反推，看到的恒星位置就比真实位置“外移”了一点——日全食时太阳背后的星星仿佛被挪了位置，这正是 1919 年爱丁顿要找的位移。
\end{mainline}

\begin{center}
\begin{tikzpicture}[scale=1.0,>=Stealth]
  % 太阳
  \fill[orange!60] (0,0) circle (0.85);
  \node at (0,0) {太阳};
  % 真实光路（折线示意偏折）
  \draw[thick] (-6,1.05) -- (0,0.95) -- (6,0.55);
  \fill (-6,1.05) circle (1.5pt);
  \node[above] at (-6,1.15) {恒星（真实位置）};
  % 视线反推的虚像方向
  \draw[dashed] (6,0.55) -- (0,1.35);
  \draw[dashed] (0,1.35) -- (-2.6,1.57);
  \fill[white,draw] (-2.6,1.57) circle (1.5pt);
  \node[above] at (-2.6,1.7) {看到的像};
  % 偏折角标注
  \draw[<->] (3.4,0.955) arc[start angle=95,end angle=86,radius=3.4];
  \node[right] at (3.9,1.15) {\(1.75''\)（示意，已夸大）};
  % 地球
  \fill[blue!50] (6,0.55) circle (2.5pt);
  \node[below] at (6,0.3) {地球};
\end{tikzpicture}
\end{center}

\begin{mathbridge}
\textbf{一、用多普勒效应推出钟差公式。} 电梯舱以加速度 \(g\) 上升，舱高 \(h\)。地板发光时舱速为零，光用约 \(t=h/c\) 的时间到达天花板，此时天花板速度已达 \(v=gt=gh/c\)。运动接收者测到的频率近似按多普勒因子压低：\(\Delta f/f\approx v/c=gh/c^{2}\)。等效原理把它翻译成引力场结论：高度差 \(h\) 处钟差 \(\Delta f/f=gh/c^{2}\)，推广到一般弱场就是 \(\Delta\Phi/c^{2}\)。整个推导只用了匀速运动学和多普勒直觉——这正是等效原理的威力：它把“引力问题”兑换成“运动学问题”。

\textbf{二、史瓦西半径与一个诚实的巧合。} 牛顿力学里，天体表面的逃逸速度是 \(v=\sqrt{2GM/r}\)。若追问“半径多小，逃逸速度会达到光速”，解出
\[
r_{s}=\frac{2GM}{c^{2}}
\]
代入真实数据：太阳的 \(r_{s}\) 约 3 千米，地球约 9 毫米，一座大山不到一个原子核。这个 \(r_{s}\) 与广义相对论严格解出的黑洞视界半径数值完全相同——但请听清楚：这只是数值上的巧合。光不是“速度不够逃不出去”的抛体，黑洞也不是“引力强到拉住光的星球”；真正的图像是视界以内时空的几何让所有的“前方”都指向中心。公式可背，图像要换。
\end{mathbridge}

\begin{challenge}
广义相对论的心脏是爱因斯坦场方程，这里只展示它的长相，不作任何推导：
\[
G_{\mu\nu}=\frac{8\pi G}{c^{4}}\,T_{\mu\nu}
\]
左边 \(G_{\mu\nu}\) 是一组描述时空弯曲程度的量（由度规及其变化率组合而成），右边 \(T_{\mu\nu}\) 描述物质和能量的分布——物质怎么摆，时空就怎么弯。它外表只有一行，展开是十个互相缠绕的偏微分方程，人类至今只在少数高度对称的情形下求得精确解：球对称恒星外部（史瓦西解）、旋转黑洞（克尔解）、均匀膨胀的宇宙（弗里德曼解）等。自由物体的运动则由测地线方程给出：在弯曲时空里走“最直的路”。

这个理论最诚实的部分是它的自我警告。方程允许“奇点”存在：黑洞中心处曲率无限大，方程自己算不下去了——这不是宇宙的真相，而是理论在举手投降，声明自己的适用范围到头了。要把奇点问到底，需要把广义相对论和第 42 章开始的量子理论缝成同一套语言，而这门“量子引力”至今没有完工：黑洞的熵、霍金辐射、宇宙最初的时刻，都卡在这条接缝上。本章的邀请函送到这里为止——门后的工地还在施工。
\end{challenge}

\step{模型的边界}

\begin{mainline}
先给等效原理划界。它说的是“局部”不可区分，不是“处处”不可区分。把电梯做得足够大、实验做得足够久，破绽就来了：在真实引力场中自由下落的舱里，并排悬浮的两个小球会慢慢互相靠近——因为它们都沿直线坠向地心，两条路线并不平行；上下叠放的两个小球会慢慢远离——因为低处的引力更强。这类因引力不均匀产生的效应叫潮汐效应（海水的潮汐正是月亮对地球两侧的“差别拉扯”）。直觉会说“失重环境里物体该静静悬浮，纹丝不动”，这两个小球就是反例：失重是真的，相对漂移也是真的。所以准确的说法是：在一个足够小的区域、足够短的时间里，引力等效于加速；一旦尺度过大，潮汐会把引力“供出来”。

再给弱场公式划界。\(\Delta f/f=\Delta\Phi/c^{2}\) 在太阳系内好用到不可思议，但在中子星表面、黑洞视界附近，它只配当粗略估计。广义相对论自身也有边界：奇点处失效，与量子力学尚未统一。这不叫“错了”，而叫适用范围的尽头——正如牛顿力学是它在低速弱场下的近似，广义相对论也等着被更深的理论包容。

接下来处理一个必须当心的比喻：蹦床。流行科普里，太阳被画成蹦床上的重球，把布面压出凹陷，地球沿着凹陷滚圈。这个图像像什么？它像的地方只有一条：几何形状确实会引导运动方向，正如飞机沿大圆飞行。它不像什么？它不像的地方太多了：布面是被地球引力压弯的，用它“解释引力”是拿引力解释引力的循环论证；布是二维的，真实弯曲的是四维时空，而且弯曲的主角是时间而非空间；球在布上会因摩擦越滚越慢，行星在弯曲时空里却永不疲倦。请把这个比喻当作一张门口的照片，而不是房间本身。

还有一个会绊倒直觉的反例必须说清楚：黑洞不是宇宙吸尘器。如果太阳此刻塌缩成同质量的黑洞，地球的轨道几乎纹丝不动——决定轨道的是质量如何弯曲外部时空，质量没变，外部几何就没变。你不会被“吸”进去，除非你本来就走得太近：靠到视界附近，时空几何才会露出獠牙。黑洞的可怕不在“吸力大”，在于它有一个只进不出的边界。

最后按全书规矩分清三层。实验事实：自由落体的普适性（\(10^{-15}\) 精度）、庞德—雷布卡塔实验、飞行原子钟、GPS 的 38 微秒、1919 年以来的光线偏折、水星近日点、2015 年直接探测到的引力波、近年来银河系中心黑洞的成像——这些谁去测都一样。数学模型：等效原理、度规几何、场方程、测地线——它们把事实组织成可以计算的规则，弱场低速时退化为牛顿引力。物理解释：“引力就是时空几何”是目前最好用的读法，但它是对模型的诠释；未来更深的理论也许会换一个更底层的说法，正如广义相对论换掉了“万有引力”的说法。把三层分开，你既不会轻视这个理论，也不会把它神化。
\end{mainline}

\step{回头解释开场现象}

\begin{mainline}
回到开场那两笔账。

第一笔：卫星钟为什么每天快 38 微秒？现在它可以拆成两项写进账簿。广义相对论项：卫星在两万千米高处，引力势比地面“浅”，钟天然走得快，每天快约 46 微秒；狭义相对论项：卫星以每秒 3.9 千米飞奔，运动使它慢，每天慢约 7 微秒。合计每天净快约 38 微秒，乘以光速得每天约 11 千米的定位漂移——与工程部门实测和补偿的数目一致。卫星出厂时，原子钟被预先调慢约每天 38 微秒，上天后两种相对论效应刚好把节拍“拨回”地面标准。你的每一次导航，都是广义相对论的一次日常签到。

第二笔：宇航员为什么失重？不是引力消失了——四百千米高空的引力仍有地面的九成。真相是：空间站和宇航员都在沿同一条弯曲时空里的“最直的路”自由下落，只不过这条路被地球弯成了一圈又一圈的椭圆。你身边所有物体和你走同一条路，彼此之间没有支撑、没有拉扯，于是一切飘浮。你在动手实验里抛起手机的那零点几秒，传感器读数归零，你和一枚微型空间站毫无区别。所谓“失重”，准确的名字叫“自由下落”；所谓“体重”，是地面不许你自由下落时给你的推力。

顺带兑现本章标题里的另外两张请柬。引力波：当大质量天体剧烈运动——比如两个黑洞相互绕转、最终并合——时空的弯曲本身会像涟漪一样以光速扩散。2015 年 9 月，美国的激光干涉引力波天文台首次直接听到这种涟漪：两束互相垂直的四千米长光臂，被扫过的引力波交替拉伸压缩，臂长变化只有 \(10^{-18}\) 米量级，约是质子直径的千分之一，却被激光干涉稳稳读出。波形与一对三十倍太阳质量上下的黑洞并合的广义相对论预言逐点吻合，并合中约有三个太阳质量的能量以引力波形式泼洒出去。引力透镜：星系团的庞大质量把时空压成一块“透镜”，背后星系的光被弯成弧、环甚至多重像——天文学家反过来用这些弧线给看不见的质量“称重”，绘出暗物质的分布图，还用恒星级质量的微透镜捕捉到了一批系外行星。引力不再是拉动物的绳子，而成了观测宇宙的望远镜。

最后看一眼引力最极端的舞台。银河系的正中心蹲着一个约四百万倍太阳质量的致密天体。天文学家花二十多年追踪它周围的恒星，其中一颗编号 S2 的星以十六年为期绕着它狂奔，近日点速度高达光速的百分之二点几，它在近日点发出的光被测到了与广义相对论预言相符的引力红移。2019 年和 2022 年，事件视界望远镜用横跨地球的射电干涉阵，先后拍下了星系 M87 中心和银河系中心这两个黑洞的“剪影”：发光的炽热气体环绕着一团暗影，暗影的大小与理论计算的视界附近结构一致。一百年前从电梯里想出来的理论，如今有了可以贴在墙上的照片。
\end{mainline}

\step{脑洞实验与挑战题}

\begin{thought}[把地球塞进一颗葡萄，再站上中子星]
做两个极端尺度的推算。第一个：把地球整个压缩，压到半径 9 毫米——约一粒葡萄的大小——逃逸速度的牛顿式估算撞上光速，它就成了黑洞。质量一点没变，如果此时月亮还在原距离，它将照常绕这颗“葡萄”转圈，分毫不变。黑洞的怪物属性不来自质量，来自“把质量塞进极小体积”之后时空几何的剧烈扭曲。

第二个：中子星是真实存在的天体，约 1.4 倍太阳质量，半径只有约 10 千米——把比太阳还重的物质压进一座城市。它的史瓦西半径约 4 千米，与半径之比约 0.4，表面的钟比远处慢约两成：那里过 80 年，我们这里过 100 年。如果你在中子星表面站一秒钟，你脚下的引力约为地面的一千亿倍——当然你根本站不住，潮汐力会先把任何物体沿径向拉长。这颗“城市大小的原子核”不是科幻，射电望远镜日复一日地收到它们灯塔般的脉冲，周期稳定到可以当钟用，毫秒级脉冲星的计时阵列甚至被用来聆听整个星系尺度的引力波背景。

最后留一个安静的问题：你的头和你的脚，离地心差了约一米七，按钟差公式，你的脚比你的头老得慢——一辈子下来相差约几十纳秒。这个差别永远影响不了你的人生，但它是真实的：时间在你身上都不是均匀的。“此刻”这个词，在引力世界里从来就不是宇宙的公共财产。
\end{thought}

\begin{puzzles}
\item 电梯里的一束光。在地面静止的电梯里，从舱壁水平射出一束激光，舱高 2 米。估算光打到对面舱壁时比入射点低多少？为什么日常生活中从来没人看见过手电的光往下弯？（思路提示：把舱当成以 \(g\) 加速，光飞行时间约为舱宽除以 \(c\)，哪怕舱宽 3 米，时间也只有 \(10^{-8}\) 秒量级，“下落”距离是 \(gt^{2}/2\sim10^{-15}\) 米量级，比原子还小。引力弯光是真实的，但地球的 \(g\) 配上光速，效应小到荒谬——所以天文学家要借太阳这么大的质量、日食这么好的舞台才看得见。）
\item 两只钟的约定。一对双胞胎，甲留在海平面，乙搬到海拔四千米的山顶住五十年，重逢时谁更老？大约老多少？这个差别和“谁运动得快”有没有关系？（思路提示：用 \(\Delta f/f=gh/c^{2}\)，\(h=4000\) 米时 \(\Delta f/f\approx4\times10^{-13}\)，五十年约快 0.6 毫秒，乙更老。这是引力势的差别，与谁爬过山、谁坐车无关；第 39 章那种“运动钟慢”的效应在这里是另一笔小得多的账。）
\item 一半是时间，一半是空间。用等效原理的电梯粗估星光掠过太阳的偏折，只得 0.87 角秒，而实测是 1.75 角秒。这个“因子 2”说明了什么？为什么说它既是等效原理的胜利、又是等效原理的边界？（思路提示：电梯推理只用了“钟在引力场中变慢”，即时间的弯曲；完整理论里空间本身的几何也弯曲，两份贡献相加才对。胜利在于：直觉不用任何场方程就拿到了正确的量级和正确的物理机制；边界在于：局部直觉无法替代完整的时空几何，精确预言必须回到场方程。）
\item 听不见的涟漪为什么难测。引力波扫过地球时，把四千米长的干涉臂改变了约 \(10^{-18}\) 米。有人质疑：“这么小的变化，肯定是仪器噪声，不算发现。”请从“相对变化”和“多台探测器互相印证”的角度，说说为什么科学家敢宣布这是真实的信号。（思路提示：引力波改变的是长度本身的比例，约 \(10^{-21}\) 的应变对四千米臂长才是 \(10^{-18}\) 米——它不是尺子量不出，而是需要干涉仪把光程差放大到可测；更关键的是，两座相距三千千米的探测器在 7 毫秒内先后记录到形状一致的波形，噪声无法跨大陆合谋，信号还以光速在两站之间传播，方向也对得上。）
\end{puzzles}
